% ============================================================
% TUTORIAL: OCI Generative AI — De la configuración al primer chat
% Proyecto: Sentinel AcademIA
% Autor: Setup Log (lFunknown)
% Compilar: tectonic tutorial-oci-llm.tex
% ============================================================

\documentclass[11pt,a4paper,oneside]{article}

% --- Paquetes base ---
\usepackage[utf8]{inputenc}
\usepackage[T1]{fontenc}
\usepackage[spanish,es-nodecimaldot]{babel}
\usepackage{lmodern}
\usepackage[margin=2.5cm]{geometry}
\usepackage{parskip}
\usepackage{microtype}

% --- Tablas, colores, cajas ---
\usepackage{xcolor}
\usepackage{colortbl}
\usepackage{array}
\usepackage{booktabs}
\usepackage{tcolorbox}
\tcbuselibrary{breakable,skins}

% --- Código fuente con syntax highlighting ---
\usepackage{listings}

% Definir lenguaje JSON (no viene por defecto)
\lstdefinelanguage{json}{
  basicstyle=\ttfamily\small,
  showstringspaces=false,
  breaklines=true,
  literate=
    *{\{}{{{\{}}}{1}
     {\}}{{{\}}}}{1}
     {[}{{{[}}}{1}
     {]}{{{]}}}{1}
     {:}{{{:}}}{1}
     {,}{{{,}}}{1},
  string=[s]{"}{"},
  stringstyle=\color{red!70!black},
  morestring=[b]",
  morecomment=[l]{:},
  commentstyle=\color{green!50!black}
}

\lstset{
  basicstyle=\ttfamily\small,
  keywordstyle=\color{blue!70!black}\bfseries,
  stringstyle=\color{red!70!black},
  commentstyle=\color{green!50!black}\itshape,
  numberstyle=\tiny\color{gray},
  numbers=left,
  numbersep=8pt,
  frame=single,
  framerule=0.4pt,
  rulecolor=\color{gray!50},
  backgroundcolor=\color{gray!5},
  breaklines=true,
  breakatwhitespace=true,
  showstringspaces=false,
  captionpos=b,
  tabsize=2,
  literate={á}{{\'a}}1 {é}{{\'e}}1 {í}{{\'i}}1 {ó}{{\'o}}1 {ú}{{\'u}}1
           {ñ}{{\~n}}1 {Á}{{\'A}}1 {É}{{\'E}}1 {Í}{{\'I}}1 {Ó}{{\'O}}1 {Ú}{{\'U}}1
           {¿}{{?`}}1 {¡}{{!`}}1
}

% --- Hipervínculos ---
\usepackage[colorlinks=true,linkcolor=blue!60!black,urlcolor=blue!60!black,citecolor=blue!60!black]{hyperref}

% --- Encabezados ---
\usepackage{fancyhdr}
\pagestyle{fancy}
\fancyhf{}
\fancyhead[L]{\small\textsc{Sentinel AcademIA}}
\fancyhead[R]{\small Tutorial OCI LLM}
\fancyfoot[C]{\thepage}
\renewcommand{\headrulewidth}{0.4pt}

% --- Colores personalizados ---
\definecolor{okgreen}{RGB}{40,140,60}
\definecolor{errred}{RGB}{200,50,50}
\definecolor{warnorange}{RGB}{220,130,30}
\definecolor{infoblue}{RGB}{30,90,180}
\definecolor{codebg}{RGB}{245,245,245}

% --- Comandos personalizados ---
\newcommand{\cmd}[1]{\texttt{\small #1}}
\newcommand{\filename}[1]{\texttt{\small #1}}
\newcommand{\ociurl}{https://docs.oracle.com/en-us/iaas/Content/generative-ai/}
\newcommand{\ok}{\textcolor{okgreen}{\textbf{\checkmark}}}
\newcommand{\err}{\textcolor{errred}{\textbf{\ding{55}}}}
\newcommand{\warn}{\textcolor{warnorange}{\textbf{\textasciitilde}}}

% --- Metadatos ---
\title{\Huge\textbf{OCI Generative AI} \\[0.3em]
       \Large De la configuración al primer chat \\[0.5em]
       \normalsize Tutorial práctico paso a paso}
\author{Proyecto Sentinel AcademIA \\ \small lFunknown}
\date{\today}

% ============================================================
\begin{document}

\maketitle
\thispagestyle{empty}

\begin{tcolorbox}[colback=blue!5!white,colframe=infoblue,title={\textbf{¿Qué aprenderás en este tutorial?}}]
Este documento es el \textbf{log de aprendizaje} completo de cómo configurar y usar OCI Generative AI Inference desde tu terminal Linux. Está pensado para que lo leas y lo practiques tú mismo, replicando los pasos en tu propia cuenta.

\textbf{Lo que cubrimos:}
\begin{itemize}
  \item Por qué el estándar moderno es \textbf{chat} (no \texttt{text})
  \item Cómo \textbf{Cada modelo} (Cohere vs Gemini) tiene un formato distinto
  \item La diferencia entre \textbf{OCID del modelo} y \textbf{alias} (y por qué importa)
  \item Errores reales que encontramos y cómo resolverlos
  \item SDK Python \textbf{y} CLI: diferencias, cuándo usar cada uno
  \item Ejercicios prácticos al final
\end{itemize}
\end{tcolorbox}

\begin{tcolorbox}[colback=warnorange!5!white,colframe=warnorange,title={\textbf{Disclaimer: Oracle + IA cambian constantemente}}]
OCI Generative AI es un servicio \textbf{en evolución rápida}. Lo que era verdad en enero puede no serlo en junio:

\begin{itemize}
  \item Modelos nuevos aparecen cada mes (Llama 4, Grok 4, Gemini 2.5\ldots)
  \item Alias como \texttt{cohere.command-latest} pueden \textbf{redirigir} a otro modelo
  \item APIs deprecadas (\texttt{generate-text} $\rightarrow$ \texttt{chat})
  \item Regiones se expanden, OCIDs cambian de formato
\end{itemize}

\textbf{Lección:} para producción, prefiere \textbf{OCID} sobre \textbf{alias}. Para aprender, los alias est\'an bien.
\end{tcolorbox}

\vspace{1em}

% ============================================================
\section{Introducción: el cambio de paradigma \texttt{generate-text} $\rightarrow$ \texttt{chat}}
% ============================================================

OCI Generative AI ofrece dos grandes categorías de operaciones:

\begin{center}
\begin{tabular}{@{}lll@{}}
\toprule
\textbf{Operación} & \textbf{Sustituye} & \textbf{Cuándo usar} \\
\midrule
\rowcolor{gray!10}
\texttt{generate-text} & (legacy) & Modelos viejos, deprecated \\
\texttt{chat} & \textbf{ESTÁNDAR} & \textbf{Todo modelo nuevo} \\
\bottomrule
\end{tabular}
\end{center}

\subsection{¿Por qué \texttt{chat} y no \texttt{text}?}

\begin{itemize}
  \item \textbf{Historial}: \texttt{chat} mantiene contexto conversacional (lista de mensajes con roles)
  \item \textbf{Multimodalidad}: \texttt{chat} acepta imágenes, archivos, audio (en modelos que lo soportan)
  \item \textbf{Tool calling}: \texttt{chat} permite que el modelo invoque funciones externas
  \item \textbf{RAG nativo}: \texttt{chat} integra búsqueda en bases de conocimiento Oracle
  \item \textbf{El SDK deprecation}: cualquier intento de \texttt{generate-text} con modelos nuevos retorna:
\end{itemize}

\begin{tcolorbox}[colback=errred!5!white,colframe=errred]
\textbf{Error típico:} \\
\texttt{Model cohere.command-latest does not support TextGeneration, only Chat.}
\end{tcolorbox}

\subsection{El cambio de mentalidad: de ``prompt'' a ``messages''}

El formato \texttt{text} espera un solo string:

\begin{lstlisting}[language=Python,caption={generate-text (obsoleto)}]
prompt = "Resume este articulo en una oracion"
response = client.generate_text(
    inference_request=inference_request_with_prompt=prompt
)
\end{lstlisting}

El formato \texttt{chat} espera una lista de mensajes con roles:

\begin{lstlisting}[language=Python,caption={chat (estándar actual)}]
messages = [
    {"role": "SYSTEM", "content": [{"type": "TEXT", "text": "Eres un asistente conciso."}]},
    {"role": "USER",   "content": [{"type": "TEXT", "text": "Resume este articulo"}]},
]
response = client.chat(chat_request=CohereChatRequest(message="..."))
\end{lstlisting}

\textbf{Conclusión:} en este tutorial \textbf{siempre} usaremos \texttt{chat}.

% ============================================================
\section{Requisitos previos}
% ============================================================

\subsection{Herramientas a instalar}

\begin{center}
\begin{tabular}{@{}lll@{}}
\toprule
\textbf{Herramienta} & \textbf{Versión mínima} & \textbf{Para qué} \\
\midrule
\texttt{oci} (CLI)        & 3.0+  & Llamadas desde terminal \\
Python                    & 3.10+ & SDK de OCI \\
\texttt{oci} SDK (pip)    & 2.130+& Llamadas programáticas \\
\bottomrule
\end{tabular}
\end{center}

\subsubsection{Instalación de OCI CLI (vía instalador oficial de Oracle)}

El instalador de Oracle empaqueta Python 3.14 \textbf{y} la CLI juntas en un solo lugar (\texttt{\textasciitilde/lib/oracle-cli/bin/}). Esto evita contaminar tu Python del sistema.

\begin{lstlisting}[language=bash,caption={Instalar OCI CLI en Linux}]
# Descargar de https://www.oracle.com/database/technologies/instant-client/linux-x86-64-downloads.html
# O usar pip del sistema:
pip3 install oci-cli

# Verificar
oci --version
\end{lstlisting}

\subsubsection{Instalación del SDK de OCI para Python}

\begin{lstlisting}[language=bash,caption={SDK de Python}]
# Si usas el Python del sistema:
pip3 install oci

# Si usas el Python del instalador de Oracle (recomendado para aislar):
~/lib/oracle-cli/bin/pip3 install oci
\end{lstlisting}

\subsection{Configurar el archivo \filename{\textasciitilde/.oci/config}}

Este archivo contiene tus credenciales. \textbf{Nunca lo commitees a git.}

\begin{lstlisting}[language=bash,caption={~/.oci/config (formato INI)}]
[DEFAULT]
user=ocid1.user.oc1..aaaaaa...       # OCID del usuario IAM
fingerprint=33:c3:5e:8d:...         # Huella de la API key
tenancy=ocid1.tenancy.oc1..aaaaaa... # OCID del tenancy (root)
region=us-chicago-1                  # Región (debe coincidir con el modelo)
key_file=/home/usuario/.oci/private_key.pem
\end{lstlisting}

\textbf{Notas importantes:}

\begin{itemize}
  \item \texttt{region} \textbf{debe ser la misma} donde desplegaste el modelo (Chicago, Frankfurt, etc.)
  \item \texttt{key\_file} debe tener permisos \texttt{600}: \cmd{chmod 600 \textasciitilde/.oci/private\_key.pem} (el flag \texttt{\textasciitilde} se expande al HOME del usuario)
  \item \textbf{No} incluyas \texttt{compartment-id} en el config (lo pasamos por CLI/SDK)
  \item Si tienes varios perfiles, usa \texttt{[DEFAULT]}, \texttt{[ACADEMY]}, etc.
\end{itemize}

\subsubsection{Obtener el OCID del compartment}

\begin{lstlisting}[language=bash,caption={Listar compartments}]
oci iam compartment list --all \
  --query 'data[?name==`SentinelAcademia`].{name:name, id:id}' \
  --output table
\end{lstlisting}

El OCID que devuelve (empieza con \texttt{ocid1.compartment.oc1..}) es el que pasamos a las llamadas LLM.

% ============================================================
\section{Concepto clave 1: cada modelo tiene un ``request type'' distinto}
% ============================================================

Este es el primer punto clave del tutorial y la razón por la que te confundiste al principio.

\begin{tcolorbox}[colback=warnorange!5!white,colframe=warnorange,title={\textbf{Regla de oro \#1}}]
\textbf{El ``serving mode'' define qué tipo de objeto \texttt{chatRequest} debes usar.}
\begin{itemize}
  \item Modelo Cohere $\rightarrow$ \texttt{CohereChatRequest}
  \item Modelo Gemini $\rightarrow$ \texttt{GenericChatRequest} con contenido en formato array
  \item Mezclarlos $\rightarrow$ error 400 del API
\end{itemize}
\end{tcolorbox}

\subsection{Tabla resumen: SDK Python}

\begin{center}
\small
\begin{tabular}{@{}llll@{}}
\toprule
\textbf{Modelo (alias)} & \textbf{Request class} & \textbf{Campo de mensaje} & \textbf{Formato} \\
\midrule
\rowcolor{blue!10}
\texttt{cohere.command-latest} & \texttt{CohereChatRequest} & \texttt{message} (string) & ``Di hola'' \\
\texttt{cohere.command-r-plus}  & \texttt{CohereChatRequest} & \texttt{message} (string) & ``Di hola'' \\
\rowcolor{blue!10}
\texttt{google.gemini-2.5-flash} & \texttt{GenericChatRequest} & \texttt{messages} (lista) & lista con roles \\
\bottomrule
\end{tabular}
\end{center}

\subsection{Tabla resumen: OCI CLI}

\begin{center}
\small
\begin{tabular}{@{}lll@{}}
\toprule
\textbf{Modelo} & \textbf{Sub-comando} & \textbf{Flag del mensaje} \\
\midrule
\rowcolor{blue!10}
Cohere (cualquiera) & \texttt{chat-cohere-chat-request} & \texttt{-{}-chat-request-message} (string) \\
Gemini (cualquiera)  & \texttt{chat-generic-chat-request} & \texttt{-{}-chat-request-messages} (JSON array) \\
\bottomrule
\end{tabular}
\end{center}

\textbf{Esto NO es documentado claramente en la doc oficial de OCI.} Lo descubrimos probando.

% ============================================================
\section{Concepto clave 2: OCID del modelo vs alias (lo que m\'as nos confundi\'o)}
% ============================================================

Este es el \textbf{segundo punto clave}, igual de importante que el primero.

\subsection{Las dos formas de identificar un modelo en OCI}

\begin{center}
\begin{tabular}{@{}p{3.5cm}p{5.5cm}p{5.5cm}@{}}
\toprule
 & \textbf{Alias (modelId)} & \textbf{OCID del modelo} \\
\midrule
\rowcolor{blue!10}
\textbf{Formato} & \texttt{cohere.command-latest} & \texttt{ocid1.generativeaimodel.oc1.us-chicago-1.amaaaaaa...} \\
\textbf{Estabilidad} & Puede cambiar (Oracle lo apunta a otro modelo) & Inmutable, identifica una versi\'on exacta \\
\rowcolor{blue!10}
\textbf{Para qu\'e} & Aprender, pruebas r\'apidas, documentaci\'on & Producci\'on, deploys reproducibles \\
\textbf{D\'onde verlo} & Cualquier doc o CLI & \textbf{Solo en la Consola de OCI} \\
\bottomrule
\end{tabular}
\end{center}

\subsection{¿Por qu\'e importa?}

\textbf{Ejemplo real:} hoy \texttt{cohere.command-latest} apunta a Command R (ago-2024). Ma\~nana Oracle podr\'ia actualizarlo a Command R2 (2026) sin avisarte. Tu c\'odigo se comporta diferente sin que hayas tocado nada.

Con el OCID del modelo, eso \textbf{no pasa}: est\'as pinneado a esa versi\'on exacta para siempre.



El campo \texttt{modelId} dentro de \texttt{-{}-serving-mode} acepta \textbf{los dos} formatos:

\begin{lstlisting}[language=bash,caption={Misma llamada, dos formas de identificar el modelo}]
# Con alias
--serving-mode '{"modelId":"cohere.command-latest","servingType":"ON_DEMAND"}'

# Con OCID
--serving-mode '{"modelId":"ocid1.generativeaimodel.oc1.us-chicago-1.amaaaaaa...","servingType":"ON_DEMAND"}'
\end{lstlisting}

El servidor de OCI resuelve ambos al mismo modelo f\'isico.

% ============================================================
\section{C\'omo obtener el OCID de un modelo pre-built}
% ============================================================

\begin{tcolorbox}[colback=errred!5!white,colframe=errred,title={\textbf{Importante}}]
\textbf{La CLI de OCI no lista los modelos pre-built directamente.} Los OCIDs solo se pueden obtener desde la \textbf{Consola Web de OCI}.
\end{tcolorbox}

\subsection{Paso a paso en la Consola}

\begin{enumerate}
  \item Abre \url{https://cloud.oracle.com/}
  \item Men\'u $\rightarrow$ \textbf{Analytics \& AI} $\rightarrow$ \textbf{Generative AI} $\rightarrow$ \textbf{Models}
  \item Ver\'as una tabla con todos los modelos disponibles en tu regi\'on
  \item Cada fila tiene una columna \textbf{OCID} que empieza con \texttt{ocid1.generativeaimodel.oc1..}
  \item Copia el OCID del modelo que quieras (ej. \texttt{cohere.command-latest})
\end{enumerate}

\subsection{Verificaci\'on: ¿el OCID es correcto?}

Una vez que tienes el OCID, pr\'uebalo:

\begin{lstlisting}[language=bash,caption={Verificar que el OCID responde}]
# Reemplaza <OCID> por el que copiaste
oci generative-ai-inference chat-result \
  chat-cohere-chat-request \
  -c "$COMPARTMENT_OCID" \
  --serving-mode '{"modelId":"<OCID>","servingType":"ON_DEMAND"}' \
  --chat-request-message "Di hola" \
  --chat-request-max-tokens 20
\end{lstlisting}

Si devuelve texto $\rightarrow$ OCID v\'alido. Si devuelve 404 $\rightarrow$ revisa que sea de la misma regi\'on.

% ============================================================
\section{Concepto clave 3: estrategia de estabilidad para producci\'on}
% ============================================================

\begin{tcolorbox}[colback=okgreen!5!white,colframe=okgreen,title={\textbf{Regla de oro \#2}}]
Para \textbf{desarrollo y aprendizaje}: usa \textbf{alias} (m\'as corto, m\'as legible). \\
Para \textbf{producci\'on}: usa \textbf{OCID} (estable, reproducible).
\end{tcolorbox}

\subsection{Pattern recomendado en c\'odigo}

\begin{lstlisting}[language=Python,caption={src/services/llm\_client.py — versi\'on con OCID}]
import os
import oci
from oci.generative_ai_inference import models as m

# OCID pinneado (tomado de la Consola de OCI una sola vez)
COHERE_OCID = os.environ["COHERE_MODEL_OCID"]
COMPARTMENT_OCID = os.environ["COMPARTMENT_OCID"]

def chat(prompt: str) -> str:
    config = oci.config.from_file("~/.oci/config", "DEFAULT")
    client = oci.generative_ai_inference.GenerativeAiInferenceClient(
        config, service_endpoint="https://inference.generativeai.us-chicago-1.oci.oraclecloud.com"
    )
    details = m.ChatDetails(
        compartment_id=COMPARTMENT_OCID,
        serving_mode=m.OnDemandServingMode(model_id=COHERE_OCID),  # OCID, no alias
        chat_request=m.CohereChatRequest(message=prompt, max_tokens=500, temperature=0.3),
    )
    return client.chat(details).data.chat_response.text
\end{lstlisting}

\subsection{Alias vs OCID: cu\'ando rotar}

\begin{itemize}
  \item \textbf{Cambias OCID}: cuando \emph{t\'u} decides actualizar a un modelo nuevo (testeo, A/B)
  \item \textbf{Cambias alias}: cuando Oracle anuncia una deprecation del alias
  \item \textbf{Nunca cambias en runtime}: el OCID debe ser constante durante toda la vida del deploy
\end{itemize}

% ============================================================
\section{Guía 1: Cohere Command con Python SDK}
% ============================================================

\subsection{Código funcional (verificado) — dos versiones}

\subsubsection{Versi\'on A: con alias (para aprender)}

\begin{lstlisting}[language=Python,caption={test\_cohere\_chat.py — con alias}]
import oci
from oci.generative_ai_inference import models as m

config = oci.config.from_file("~/.oci/config", "DEFAULT")
client = oci.generative_ai_inference.GenerativeAiInferenceClient(
    config, service_endpoint="https://inference.generativeai.us-chicago-1.oci.oraclecloud.com"
)

chat_request = m.CohereChatRequest(
    message="Explica qu\'e es Python en una sola frase",
    max_tokens=100,
    temperature=0.3,
    is_stream=False,
)

chat_details = m.ChatDetails(
    compartment_id=config["tenancy"],
    serving_mode=m.OnDemandServingMode(model_id="cohere.command-latest"),  # alias
    chat_request=chat_request,
)

response = client.chat(chat_details)
print(response.data.chat_response.text)
\end{lstlisting}

\subsubsection{Versi\'on B: con OCID (para producci\'on)}

\begin{lstlisting}[language=Python,caption={test\_cohere\_chat.py — con OCID}]
import oci
import os
from oci.generative_ai_inference import models as m

# OCID copiado UNA VEZ desde la Consola de OCI
COHERE_OCID = os.environ["COHERE_MODEL_OCID"]

config = oci.config.from_file("~/.oci/config", "DEFAULT")
client = oci.generative_ai_inference.GenerativeAiInferenceClient(
    config, service_endpoint="https://inference.generativeai.us-chicago-1.oci.oraclecloud.com"
)

chat_request = m.CohereChatRequest(
    message="Explica qu\'e es Python en una sola frase",
    max_tokens=100,
    temperature=0.3,
    is_stream=False,
)

chat_details = m.ChatDetails(
    compartment_id=config["tenancy"],
    serving_mode=m.OnDemandServingMode(model_id=COHERE_OCID),  # OCID
    chat_request=chat_request,
)

response = client.chat(chat_details)
print(response.data.chat_response.text)
\end{lstlisting}

\subsection{Anatomía de la respuesta}

\begin{lstlisting}[language=Python,caption={Inspeccionar la respuesta completa}]
response.data  # -> ChatResponse
#   .chat_response -> CohereChatResponse
#       .text              -> "Python es un lenguaje..."
#       .finish_reason     -> "COMPLETE" | "MAX_TOKENS" | "ERROR"
#       .usage             -> {prompt_tokens, completion_tokens, total_tokens}
#   .model_id         -> "cohere.command-latest"
#   .model_version    -> "1.7"
\end{lstlisting}

\subsection{Errores comunes con Cohere}

\begin{tcolorbox}[colback=errred!5!white,colframe=errred,breakable]
\textbf{Error 1:} ``\texttt{Unrecognized keyword arguments: text}'' \\
\textbf{Causa:} Usaste \texttt{CohereChatRequest(text="...")} en vez de \texttt{message="..."}. \\
\textbf{Solución:} Cohere v1 usa \texttt{message} (singular, string), no \texttt{text}.

\vspace{0.5em}
\textbf{Error 2:} ``\texttt{Model does not support TextGeneration, only Chat}'' \\
\textbf{Causa:} Intentaste usar \texttt{client.generate\_text()} con un modelo Cohere nuevo. \\
\textbf{Solución:} Usa \texttt{client.chat()} siempre.

\vspace{0.5em}
\textbf{Error 3:} ``\texttt{Compartment ID must be provided}'' \\
\textbf{Causa:} Pasaste \texttt{compartment\_id} como string vacío o \texttt{None}. \\
\textbf{Solución:} Usa el OCID completo: \texttt{ocid1.compartment.oc1..aaaaaa...}
\end{tcolorbox}

% ============================================================
\section{Guía 2: Cohere Command con CLI de OCI}
% ============================================================

\subsection{Comando funcional (verificado) — dos versiones}

\subsubsection{Versi\'on A: con alias}

\begin{lstlisting}[language=bash,caption={Test Cohere CLI con alias}]
COMPARTMENT_OCID="ocid1.compartment.oc1..aaaaaaaa4vvf5a7lznbrfxporgxgmkr3mdanyktzsaosbwqc2urrq4pwwmzq"

oci generative-ai-inference chat-result \
  chat-cohere-chat-request \
  -c "$COMPARTMENT_OCID" \
  --serving-mode '{"modelId":"cohere.command-latest","servingType":"ON_DEMAND"}' \
  --chat-request-message "Di hola en una sola frase" \
  --chat-request-max-tokens 30 \
  --chat-request-temperature 0
\end{lstlisting}

\subsubsection{Versi\'on B: con OCID (recomendado para producci\'on)}

\begin{lstlisting}[language=bash,caption={Test Cohere CLI con OCID}]
# Variables de entorno (poner en ~/.bashrc)
export COMPARTMENT_OCID="ocid1.compartment.oc1..aaaaaaaa4vvf5a7lznbrfxporgxgmkr3mdanyktzsaosbwqc2urrq4pwwmzq"
export COHERE_MODEL_OCID="ocid1.generativeaimodel.oc1.us-chicago-1.amaaaaaaXXXXX..."  # <- de la Consola

oci generative-ai-inference chat-result \
  chat-cohere-chat-request \
  -c "$COMPARTMENT_OCID" \
  --serving-mode "{\"modelId\":\"$COHERE_MODEL_OCID\",\"servingType\":\"ON_DEMAND\"}" \
  --chat-request-message "Di hola en una sola frase" \
  --chat-request-max-tokens 30 \
  --chat-request-temperature 0
\end{lstlisting}

\subsubsection{Atajo: alias de shell}

Para no escribir todo el tiempo, crea este alias en tu \texttt{\textasciitilde/.bashrc}:

\begin{lstlisting}[language=bash,caption={\textasciitilde/.bashrc}]
# Requiere: COMPARTMENT_OCID y COHERE_MODEL_OCID ya exportadas
alias ocichat='oci generative-ai-inference chat-result chat-cohere-chat-request \
  -c "$COMPARTMENT_OCID" \
  --serving-mode "{\"modelId\":\"$COHERE_MODEL_OCID\",\"servingType\":\"ON_DEMAND\"}"'

# Uso:
ocichat --chat-request-message "Di hola" --chat-request-max-tokens 30
\end{lstlisting}

\subsection{Output esperado}

\begin{lstlisting}[language=json,caption={Respuesta del CLI}]
{
  "data": {
    "chat-response": {
      "api-format": "COHERE",
      "text": "¡Hola, cómo estás?",
      "finish-reason": "COMPLETE",
      "usage": {
        "prompt-tokens": 7,
        "completion-tokens": 8,
        "total-tokens": 15
      }
    },
    "model-id": "cohere.command-latest",
    "model-version": "1.7"
  }
}
\end{lstlisting}

\subsection{Anatomía de los flags}

\begin{center}
\small
\begin{tabular}{@{}lp{9cm}@{}}
\toprule
\textbf{Flag} & \textbf{Qué hace} \\
\midrule
\rowcolor{gray!10}
\texttt{chat-result} & Sub-comando de alto nivel (imprime resultado, no job) \\
\texttt{chat-cohere-chat-request} & Especifica el formato del request (Cohere v1) \\
\rowcolor{gray!10}
\texttt{-{}-compartment-id} & OCID del compartment (no del tenancy) \\
\texttt{-{}-serving-mode} & JSON con \texttt{modelId} y \texttt{servingType} \\
\rowcolor{gray!10}
\texttt{-{}-chat-request-message} & \textbf{String} con el prompt (Cohere es simple) \\
\texttt{-{}-chat-request-max-tokens} & Límite de tokens en la respuesta \\
\bottomrule
\end{tabular}
\end{center}

% ============================================================
\section{Guía 3: Gemini con CLI de OCI}
% ============================================================

\subsection{Comando funcional (verificado)}

\begin{lstlisting}[language=bash,caption={Test Gemini desde terminal}]
COMPARTMENT_OCID="ocid1.compartment.oc1..aaaaaaaa4vvf5a7lznbrfxporgxgmkr3mdanyktzsaosbwqc2urrq4pwwmzq"

oci generative-ai-inference chat-result \
  chat-generic-chat-request \
  --compartment-id "$COMPARTMENT_OCID" \
  --serving-mode '{"modelId":"google.gemini-2.5-flash","servingType":"ON_DEMAND"}' \
  --chat-request-messages '[{"role":"USER","content":[{"type":"TEXT","text":"Di hola"}]}]' \
  --chat-request-max-tokens 20
\end{lstlisting}

\subsection{¡La trampa del formato!}

\begin{tcolorbox}[colback=warnorange!5!white,colframe=warnorange,title={Diferencia clave Cohere vs Gemini}]
\textbf{Cohere CLI:} \texttt{-{}-chat-request-message} es un \textbf{string} simple.

\textbf{Gemini CLI:} \texttt{-{}-chat-request-messages} es un \textbf{JSON array} con estructura anidada:

\begin{lstlisting}[language=json]
[
  {
    "role": "USER",
    "content": [
      {
        "type": "TEXT",
        "text": "Di hola"
      }
    ]
  }
]
\end{lstlisting}

Si pasas \texttt{"content": "Di hola"} (string plano) $\rightarrow$ 400 ``Please pass in correct format of request''.
\end{tcolorbox}

\subsection{¿Por qué Gemini usa este formato raro?}

\begin{itemize}
  \item El campo \texttt{content} es una \textbf{lista} porque Gemini soporta \textbf{multimodalidad}
  \item Cada elemento de la lista es un \textbf{ContentPart} con \texttt{type}
  \item Los tipos posibles son: \texttt{TEXT}, \texttt{IMAGE}, \texttt{AUDIO}, etc.
  \item Hoy solo usamos \texttt{TEXT}, pero mañana puedes agregar imágenes
\end{itemize}

% ============================================================
\section{Comparación lado a lado: SDK vs CLI}
% ============================================================

\subsection{Cohere Command}

\begin{center}
\begin{tabular}{@{}p{6cm}p{8cm}@{}}
\toprule
\textbf{Python SDK} & \textbf{CLI} \\
\midrule
\rowcolor{codebg}
\texttt{m.CohereChatRequest(message="...")} & \texttt{-{}-chat-request-message "..."} \\
\texttt{client.chat(chat\_details)} & \texttt{oci ... chat-result} \\
\texttt{response.data.chat\_response.text} & Filtrar JSON con \texttt{jq} \\
\bottomrule
\end{tabular}
\end{center}

\subsection{Gemini}

\begin{center}
\begin{tabular}{@{}p{6cm}p{8cm}@{}}
\toprule
\textbf{Python SDK} & \textbf{CLI} \\
\midrule
\rowcolor{codebg}
\texttt{m.GenericChatRequest(messages=[...])} & \texttt{-{}-chat-request-messages '[...]'} \\
\texttt{m.TextContent(text="...")} & \texttt{"content":[ \{"type":"TEXT","text":"..."\} ]} \\
\texttt{client.chat(chat\_details)} & \texttt{oci ... chat-result} \\
\bottomrule
\end{tabular}
\end{center}

% ============================================================
\section{Troubleshooting: errores que encontramos y cómo los resolvimos}
% ============================================================

\subsection{Catálogo completo de errores}

\begin{tcolorbox}[colback=errred!5!white,colframe=errred,breakable,title={Error 1: ``No such file or directory: private\_key.pem''}]
\textbf{Cuándo:} \texttt{oci config validate} o primera llamada. \\
\textbf{Causa:} El \texttt{key\_file} en el config (ruta \texttt{\textasciitilde/.oci/config}) apunta a una ruta inexistente. \\
\textbf{Solución:}
\begin{lstlisting}[language=bash]
ls -la ~/.oci/private_key.pem
# Si no existe, descargar de la consola de OCI:
# Profile -> API Keys -> Add API Key -> Download Private Key
chmod 600 ~/.oci/private_key.pem
\end{lstlisting}
\end{tcolorbox}

\begin{tcolorbox}[colback=errred!5!white,colframe=errred,breakable,title={Error 2: ``Compartment ID is not valid''}]
\textbf{Cuándo:} Primera llamada con Python. \\
\textbf{Causa:} Pasaste el OCID entre corchetes \texttt{[ocid1...]} o con espacios. \\
\textbf{Solución:} OCID limpio, sin corchetes:
\begin{lstlisting}[language=bash]
# MAL
COMP="[ocid1.compartment.oc1..aaa]"
# BIEN
COMP="ocid1.compartment.oc1..aaa..."
\end{lstlisting}
\end{tcolorbox}

\begin{tcolorbox}[colback=errred!5!white,colframe=errred,breakable,title={Error 3: ``Entity with key cohere.command-r not found''}]
\textbf{Cuándo:} Listar modelos o primera llamada. \\
\textbf{Causa:} El nombre \texttt{cohere.command-r} sin sufijo de fecha no existe. \\
\textbf{Solución:} Usar el alias \texttt{cohere.command-latest} (siempre apunta al más nuevo) o la versión con fecha \texttt{cohere.command-r-08-2024}.
\end{tcolorbox}

\begin{tcolorbox}[colback=errred!5!white,colframe=errred,breakable,title={Error 4: ``cohere.command-latest does not support TextGeneration''}]
\textbf{Cuándo:} Usar \texttt{client.generate\_text()} con Cohere. \\
\textbf{Causa:} \texttt{generate\_text} es para modelos viejos. \\
\textbf{Solución:} Usar \texttt{client.chat()} con \texttt{CohereChatRequest} siempre.
\end{tcolorbox}

\begin{tcolorbox}[colback=errred!5!white,colframe=errred,breakable,title={Error 5: ``Unrecognized keyword arguments: text''}]
\textbf{Cuándo:} Crear \texttt{CohereChatRequest}. \\
\textbf{Causa:} Usar \texttt{text="..."} en vez de \texttt{message="..."}. \\
\textbf{Solución:} Cohere v1 usa \texttt{message} (singular).
\end{tcolorbox}

\begin{tcolorbox}[colback=errred!5!white,colframe=errred,breakable,title={Error 6: ``Please pass in correct format of request'' (Gemini CLI)}]
\textbf{Cuándo:} Llamar Gemini con \texttt{content} como string. \\
\textbf{Causa:} Gemini requiere \texttt{content} como lista de \texttt{ContentPart}. \\
\textbf{Solución:}
\begin{lstlisting}[language=bash]
# MAL
'{"role":"USER","content":"Di hola"}'

# BIEN
'{"role":"USER","content":[{"type":"TEXT","text":"Di hola"}]}'
\end{lstlisting}
\end{tcolorbox}

\begin{tcolorbox}[colback=errred!5!white,colframe=errred,breakable,title={Error 7: ``404 Not Found'' o timeout en llamada al API}]
\textbf{Cuándo:} Primera llamada desde una región nueva. \\
\textbf{Causa:} La región del \texttt{config} no coincide con la región donde está desplegado el modelo. \\
\textbf{Solución:} Verificar con:
\begin{lstlisting}[language=bash]
oci iam region list --output table
\end{lstlisting}
y usar la misma región en \texttt{-{}-region} o en el \texttt{config}.
\end{tcolorbox}

% ============================================================
\section{Ejercicios prácticos}
% ============================================================

Haz estos ejercicios en orden, \textbf{anotando} los errores que te salen y cómo los resuelves.

\subsection{Ejercicio 1: Tu primer chat con Cohere (5 min)}

\begin{enumerate}
  \item Abre tu terminal
  \item Exporta la variable \texttt{COMPARTMENT\_OCID} con tu OCID real
  \item Copia y pega el comando de la sección 5.2
  \item Verifica que el JSON de respuesta tiene \texttt{"text": "..."} con un saludo
\end{enumerate}

\textbf{Output esperado:} Un JSON con \texttt{"finish-reason": "COMPLETE"} y un texto de respuesta.

\subsection{Ejercicio 2: Cambia la temperatura (5 min)}

\begin{enumerate}
  \item Re-ejecuta el comando de Cohere con \texttt{-{}-chat-request-temperature 1.5}
  \item Compara la respuesta con la de temperatura 0
  \item ¿Qué cambió? (Creatividad, longitud, formato)
\end{enumerate}

\subsection{Ejercicio 3: Prueba Gemini (5 min)}

\begin{enumerate}
  \item Cambia el sub-comando a \texttt{chat-generic-chat-request}
  \item Cambia \texttt{-{}-chat-request-message} por \texttt{-{}-chat-request-messages} con el formato array
  \item Ejecuta y compara la velocidad con Cohere
\end{enumerate}

\subsection{Ejercicio 4: Convierte el comando CLI a Python (15 min)}

\begin{enumerate}
  \item Toma el comando CLI de Cohere que ya te funciona
  \item Trad\'{u}celo a un script Python usando \texttt{oci.generative\_ai\_inference}
  \item Guarda el script como \filename{mi\_primer\_chat.py}
  \item Ejec\'{u}talo y compara el output con el del CLI
\end{enumerate}

\textbf{Pista:} Usa el código de la secci\'on 4.1 (versi\'on A o B) como base.

\subsection{Ejercicio 5: Inserta el script en una función reutilizable (20 min)}

\begin{enumerate}
  \item Crea \filename{src/services/llm\_client.py}
  \item Define una función \texttt{chat(prompt: str) $\rightarrow$ str} que oculte toda la complejidad
  \item Escribe 3 tests con \texttt{pytest} que verifiquen:
  \begin{itemize}
    \item Devuelve string no vacío
    \item Falla limpiamente si el compartment OCID es inválido
    \item Falla limpiamente si el modelo no existe
  \end{itemize}
\end{enumerate}

\textbf{Este ejercicio es el que usaremos en Fase 1 del proyecto.}

\subsection{Ejercicio 6: Migra de alias a OCID (10 min)}

\begin{enumerate}
  \item Abre la Consola de OCI y obt\'en el OCID de \texttt{cohere.command-latest}
  \item Exporta \texttt{COHERE\_MODEL\_OCID} en tu shell
  \item Modifica tu funci\'on \texttt{chat()} para que use \texttt{COHERE\_MODEL\_OCID} en vez del alias
  \item Re-ejecuta los tests del ejercicio 5
  \item \textbf{Bonus:} crea un test que verifique que el modelo es estable (mismo OCID entre 2 llamadas)
\end{enumerate}

\textbf{Este patr\'on es el que usaremos en producci\'on.}

% ============================================================
\section{Resumen: lo que tienes que recordar}
% ============================================================

\begin{tcolorbox}[colback=okgreen!5!white,colframe=okgreen,title={\textbf{Checklist de memoria}}]
\begin{enumerate}
  \item \textbf{Est\'andar = chat}, nunca \texttt{generate\_text}
  \item \textbf{Cohere} $\rightarrow$ \texttt{CohereChatRequest} (string en \texttt{message})
  \item \textbf{Gemini} $\rightarrow$ \texttt{GenericChatRequest} (lista en \texttt{messages})
  \item \textbf{CLI Cohere} $\rightarrow$ \texttt{chat-cohere-chat-request} con string
  \item \textbf{CLI Gemini} $\rightarrow$ \texttt{chat-generic-chat-request} con array JSON
  \item \textbf{Regi\'on del config} = regi\'on del modelo
  \item \textbf{Compartment OCID} va en cada llamada, no en el config
  \item \textbf{Clave privada} con permisos 600, ubicada en \texttt{\textasciitilde/.oci/private\_key.pem}
  \item \textbf{Alias = para aprender}, \textbf{OCID = para producci\'on}
  \item \textbf{OCID del modelo solo en Consola}, no en CLI
  \item \textbf{Oracle cambia}: revisa docs cada trimestre
\end{enumerate}
\end{tcolorbox}

% ============================================================
\section{Referencias}
% ============================================================

\begin{itemize}
  \item \ociurl --- Documentación oficial de OCI Generative AI
  \item \url{https://docs.oracle.com/en-us/iaas/tools/python-sdk-examples/} --- Ejemplos del SDK de Python
  \item \url{https://www.oracle.com/artificial-intelligence/generative-ai/large-language-models/} --- Catálogo de modelos
\end{itemize}

\vspace{2em}
\begin{center}
\textit{Fin del tutorial. ¡Ahora a practicar!}
\end{center}

\end{document}
